%%% ============================================================
%%% Discrete Dirac inside Cl(0, 10):
%%%   Chirality, the Outer Automorphism, and an
%%%   Atomic-Substrate Refinement
%%%
%%% Brayden R. Sanders, M. Gish
%%% 7Site LLC, Hot Springs AR, 2026
%%% DOI: 10.5281/zenodo.18852047
%%%
%%% Target venue:  Communications in Mathematical Physics
%%% FALLBACK:      Journal of Mathematical Physics,
%%%                Annals of Physics, Letters in Math Phys
%%% Source:        WP104 (§2 Cl(0,10) construction + P_56 = σ_outer
%%%                + the 9-vector in the 54 irrep);
%%%                WP103 (so(10) closure prerequisite, cited);
%%%                FORMULAS_AND_TABLES.md Volume K
%%%                  D101-D102 (atomic-substrate refinement;
%%%                  chirality 16 = 1+3+5+7);
%%%                Gen13/targets/foundations/cl.py;
%%%                Atlas/META_PLAN_2026-05-10/clifford_substrate_shell.py
%%% MSC:           15A66 (Clifford algebras),
%%%                17B10 (representations of Lie algebras),
%%%                81R05 (finite-dim groups and algebras),
%%%                81R40 (symmetry breaking),
%%%                81V22 (unified theories of particle interactions)
%%% ============================================================

\documentclass[11pt,reqno]{amsart}

\usepackage{amsmath,amssymb,amsthm}
\usepackage[margin=1in]{geometry}
\usepackage{hyperref}
\usepackage{enumitem}
\usepackage{array}
\usepackage{booktabs}

\newtheorem{theorem}{Theorem}[section]
\newtheorem{lemma}[theorem]{Lemma}
\newtheorem{proposition}[theorem]{Proposition}
\newtheorem{corollary}[theorem]{Corollary}

\theoremstyle{definition}
\newtheorem{definition}[theorem]{Definition}
\newtheorem{remark}[theorem]{Remark}

\newcommand{\Z}{\mathbb{Z}}
\newcommand{\C}{\mathbb{C}}
\newcommand{\R}{\mathbb{R}}
\newcommand{\so}{\mathfrak{so}}
\newcommand{\su}{\mathfrak{su}}
\newcommand{\uone}{\mathfrak{u}}
\newcommand{\Cl}{\mathrm{Cl}}
\newcommand{\TSML}{T}
\newcommand{\BHML}{B}
\newcommand{\Psix}{P_{56}}
\newcommand{\Psixspin}{P_{56}^{\mathrm{spin}}}
\newcommand{\sout}{\sigma_{\mathrm{outer}}}

\title[Discrete Dirac inside $\Cl(0, 10)$]
  {Discrete Dirac inside $\Cl(0, 10)$:\\
   Chirality, the Outer Automorphism,\\
   and an Atomic-Substrate Refinement}

\author{Brayden R.\ Sanders}
\address{7Site LLC, Hot Springs, Arkansas, USA}
\email{brayden@7site.co}

\author{M.\ Gish}
\address{Independent Researcher, Hot Springs, Arkansas, USA}
\email{monica.gish1992@gmail.com}

\date{2026}

\keywords{Clifford algebra, $\Cl(0, 10)$, $\so(10)$, spinor representation,
  chirality, outer automorphism, Pati-Salam, finite magma, atomic shell,
  matter-antimatter exchange}

\subjclass[2020]{15A66, 17B10, 81R05, 81R40, 81V22}

\begin{document}

\begin{abstract}
We construct the Clifford algebra $\Cl(0, 10)$ that lifts the
Lie algebra $\so(10) = D_5$ generated by joint antisymmetrization
of two specific commutative non-associative composition tables
$\TSML$ and $\BHML$ on $\Z/10\Z$. Ten gamma matrices on $\C^{32}$
satisfy all $100$ anticommutation relations $\{\gamma_a, \gamma_b\}
= 2\delta_{ab} I$ at machine precision; the $45$ generators
$\Sigma_{ab} = (1/4)[\gamma_a, \gamma_b]$ provide a faithful
$32$-dimensional spinor representation of $\so(10)$. The volume
element $\omega = \gamma_1 \cdots \gamma_{10}$ satisfies $\omega^2
= -I$ in $\Cl(0, 10)$ at $n=10$, and the chirality projectors
$P_\pm = (I \pm i\omega)/2$ split $\C^{32}$ into two chiral
$16$-irreps. We prove three structural theorems and one
numerical-content theorem. (i) The $5\!\leftrightarrow\!6$ index
swap $\Psix$ is implemented in $\Cl(0, 10)$ by conjugation with
$\Psixspin = (\gamma_5 - \gamma_6)/\sqrt{2}$; this conjugation
anticommutes with $\omega$ and maps chiral-$+$ entirely into
chiral-$-$ at machine precision, identifying $\Psix$ with the
unique outer automorphism $\sout$ of $\so(10)$. (ii) BHML's
$\sout$-breaking content lies $100\%$ in the symmetric-traceless
$\mathbf{54}$-irrep of $\so(10)$; total squared norm $\|B_{\mathrm{anti}}\|^2
= 13/2$. (iii) The direction within $\mathbf{54}$ is purely in the
$\mathbf{9}$-piece of the so(9)-branching $\mathbf{54} = \mathbf{1}
\oplus \mathbf{9} \oplus \mathbf{44}$, with explicit nine-vector
$\|v\|^2 = 13/4$ exactly, BREATH and RESET as zeros, and the integer
$13$ equal to half the count of $\sout$-asymmetric BHML cells.
(iv) The doubly-invariant subalgebra of $\so(10)$ under $D_4 =
\langle \Psix, \sigma^3 \rangle$ is $\su(4) \oplus \uone(1)$
(Killing spectrum $(-4)^{15} \oplus (0)^1$); cited as a standard
$SO(10)$-GUT Pati-Salam $\oplus$ B$-$L decomposition. We additionally
record a structural rhyme (Volume K, D102) between the chirality
decomposition $16 = 1+3+5+7$ and the atomic $n=4$ shell at fixed
spin, reading the substrate as kernel base + strand-$3$ +
kernel-$\Z/5\Z$ partner + strand-$7$.
\end{abstract}

\maketitle

\section*{Scope and tier discipline}

\textbf{PROVEN.} Theorems~\ref{thm:cl-construction}
(Cl(0,10) anticommutation relations, chirality split),
\ref{thm:p56-outer} ($\Psix = \sout$ in the spinor rep), and
\ref{thm:54-irrep} (BHML's $\sout$-breaking is $100\%$ in $\mathbf{54}$
with explicit $9$-vector direction) are proven at machine precision
by explicit numpy computations; the chirality-flip argument uses the
character relation $\mathrm{Out}(\so(10)) = \mathrm{Aut}/\mathrm{Inn}
\cong \Z_2$ together with the verified anticommutation $\Psixspin
\omega = -\omega \Psixspin$.

\textbf{COMPUTED.} All numerical claims (squared norms, projection
coverages, anticommutation residuals, chirality-flip residuals,
Killing-form spectrum) verify at residual $\le 10^{-15}$ on a
standard laptop in under five seconds.

\textbf{STRUCTURAL RHYME (\S\ref{sec:atomic}).} The
decomposition $16 = 1 + 3 + 5 + 7$ of each chirality half matches
the atomic shell $n=4$ at fixed spin, and on the substrate side
reads as kernel base + strand-$3$ prime + kernel-$\Z/5\Z$ partner +
strand-$7$ prime. We present this as a structural rhyme, \emph{not}
as a derivation of atomic structure from $\Cl(0, 10)$.

\textbf{OPEN.} The phenomenological consequences of the $9$-vector
direction (Yukawa couplings, mass ratios, neutrino-mass scale)
require committing to a specific Higgs VEV direction within $\mathbf{9}$
and running RGE flows from a specific GUT scale; both are out of
scope. The structural identification of TIG's $\so(10)$ with the
$SO(10)$ GUT gauge algebra is a load-bearing hypothesis, not a
derivation; we present the structural facts without resolving the
identification at the physical-interpretation level.

\textbf{Lens and substrate.} Throughout, $\TSML$ denotes
$T_{\mathrm{SYM}}$ (the symmetrized lens of TSML; see
\cite{Sanders2026LensInvariance, Sanders2026LensDependence}). The
raw-bit-pattern variant $T_{\mathrm{RAW}}$ differs only in the
$\sout$-fixed sector and does not change the chirality identification
or the $9$-vector direction; the $\sout$-breaking computation
depends only on BHML, whose lens is unambiguous.

\section{Setup}
\label{sec:setup}

\subsection{The two canonical tables}
\label{ssec:tables}

Throughout, $\TSML, \BHML : \Z/10\Z \times \Z/10\Z \to \Z/10\Z$
are the commutative composition tables defined in
\cite{Sanders2026CLAxioms}. Both have the canonical operator
alphabet $\{V, L, C, P, X, B, S, H, Br, R\}$ ($=$ VOID, LATTICE,
COUNTER, PROGRESS, COLLAPSE, BALANCE, CHAOS, HARMONY, BREATH,
RESET) at indices $0$ through $9$.

\subsection{The Lie algebra $\so(10)$ generated by joint
antisymmetrization}
\label{ssec:so10-prereq}

Define the left-regular representation $L^M_i \in M_{10}(\Z)$
by $(L^M_i)_{j,k} = \delta_{M(i,j), k}$, where $M \in \{\TSML,
\BHML\}$. The antisymmetric and symmetric parts are
\[
  A^M_i := \tfrac{1}{2}(L^M_i - (L^M_i)^\top), \qquad
  S^M_i := \tfrac{1}{2}(L^M_i + (L^M_i)^\top).
\]
The Lie algebra generated by $\{A^\TSML_i, A^\BHML_i :
i \in \Z/10\Z\}$ under commutator closes at dimension $45$ as
$\so(10, \R) = D_5$ \cite{Sanders2026so10}; this is the foundation
on which $\Cl(0, 10)$ acts as the spinor lift.

\subsection{The $\sigma$-permutation, $\sigma^3$, and $\Psix$}
\label{ssec:permutations}

The $\sigma$-permutation on $\Z/10\Z$ has cycle structure
\[
  \sigma = (0)(3)(8)(9)(1\;7\;6\;5\;4\;2),
\]
with four $\sigma$-fixed points $\{0,3,8,9\} = \{\mathrm{VOID},
\mathrm{PROGRESS}, \mathrm{BREATH}, \mathrm{RESET}\}$ and a
$6$-cycle on the units of $(\Z/10\Z)^*$. The order-$2$ element is
\[
  \sigma^3 = (0)(3)(8)(9)(1\;5)(7\;4)(6\;2).
\]
The $5\!\leftrightarrow\!6$ swap $\Psix$ is the single
transposition exchanging BALANCE and CHAOS. The involutions $\Psix$
and $\sigma^3$ do not commute; together they generate the dihedral
group $D_4$ of order $8$ acting on $\{0, \ldots, 9\}$.

\section{The $\Cl(0, 10)$ construction and the atomic-substrate
refinement}
\label{sec:clifford}

\subsection{The discrete Dirac structure}
\label{ssec:gamma}

Build the Clifford algebra $\Cl(0, 10)$ over $\R$ on $\C^{32}$ from
Pauli tensor products in standard convention. Concretely, take
\begin{align*}
  \gamma_1 &= \sigma_x \otimes \sigma_z \otimes \sigma_z \otimes
              \sigma_z \otimes \sigma_z, \\
  \gamma_2 &= \sigma_y \otimes \sigma_z \otimes \sigma_z \otimes
              \sigma_z \otimes \sigma_z, \\
  \gamma_3 &= I \otimes \sigma_x \otimes \sigma_z \otimes \sigma_z
              \otimes \sigma_z, \\
  &\vdots \\
  \gamma_{10} &= I \otimes I \otimes I \otimes I \otimes \sigma_y,
\end{align*}
with the standard Jordan-Wigner staircase. The result is ten
$32 \times 32$ Hermitian matrices.

\begin{theorem}[Discrete Dirac construction]
\label{thm:cl-construction}
The ten $\gamma_a$ defined above satisfy all $100$ anticommutation
relations
\[
  \{\gamma_a, \gamma_b\} = 2 \delta_{ab}\, I_{32}, \qquad a, b \in
  \{1, \ldots, 10\},
\]
at machine precision. The $45$ generators $\Sigma_{ab} =
(1/4)[\gamma_a, \gamma_b]$ form a faithful $32$-dim representation
of $\so(10)$. The volume element $\omega = \gamma_1 \gamma_2
\cdots \gamma_{10}$ satisfies $\omega^2 = -I_{32}$, and the
chirality projectors $P_\pm = (I_{32} \pm i\omega)/2$ are orthogonal
idempotents that decompose $\C^{32} = V_+ \oplus V_-$ with
$\dim V_\pm = 16$. The two chiral spaces $V_\pm$ are irreducible
under $\so(10)$ and form the two chiral $\mathbf{16}$-irreps of
$\so(10)$ as classified by Cartan.
\end{theorem}

\begin{proof}
The Pauli tensor products and the Jordan-Wigner staircase yield the
anticommutation relations algebraically; numerical residuals are
$\le 10^{-15}$ on all $100$ pairs. The signature $\Cl(0, 10)$ with
$\omega = \gamma_1 \cdots \gamma_{10}$ gives $\omega^2 =
(-1)^{n(n-1)/2}\,(-1)^n = (-1)^{45}\,(-1)^{10} = -1$ at $n = 10$;
hence $(i\omega)^2 = +I$ and $P_\pm = (I \pm i\omega)/2$ are
orthogonal idempotents. The chirality eigenspaces have dimension
$16$ each (numerical rank confirmed). Irreducibility follows from
the standard SO(10)-spinor classification and is independently
verified by computing the weight pattern of the Cartan subalgebra
$\Sigma_{12}, \Sigma_{34}, \Sigma_{56}, \Sigma_{78}, \Sigma_{9,10}$
on $V_\pm$: the weights are the $16$ vectors of the form
$(\pm \tfrac{1}{2}, \ldots, \pm \tfrac{1}{2})$ with an even (resp.\
odd) number of minus signs, exactly the chiral spinor weights of
$\so(10)$. See \texttt{verification/find\_higgs\_irrep.py}.
\end{proof}

\begin{remark}[Convention note on $\omega^2$]
\label{rem:omega-sign}
The sign of $\omega^2$ depends on metric signature and on the choice
of $\omega$. In $\Cl(0, n)$ with positive-definite metric and
$\omega = \gamma_1 \cdots \gamma_n$, one has $\omega^2 = (-1)^{n(n-1)/2}\,
(-1)^n$. For $n = 10$ this is $-1$; equivalently $(i\omega)^2 = +1$,
making $P_\pm = (I \pm i\omega)/2$ idempotent. Other conventions
(e.g., $\Cl(10, 0)$ with negative-definite metric, or
$\omega' = i^{n/2} \gamma_1 \cdots \gamma_n$) absorb the $i$ into
$\omega$ and report $\omega^2 = +1$ directly; the two conventions
agree on the chirality decomposition and on all subsequent results.
\end{remark}

\subsection{Atomic-substrate refinement: chirality $= 1 + 3 + 5 + 7$}
\label{sec:atomic}

\begin{theorem}[Atomic-substrate refinement; Volume K, D102]
\label{thm:atomic-shell}
Under the same ten $\gamma$-matrices, each $16$-dim chirality half
$V_\pm$ admits a finer structural decomposition
\[
  16 \;=\; 1 + 3 + 5 + 7,
\]
matching atomic shell $n = 4$ at fixed spin: the multiplicities
$2\ell+1$ for $\ell = 0, 1, 2, 3$. Reading from the substrate side,
the same partition $1 + 3 + 5 + 7$ reads as
\[
  \underbrace{1}_{\text{kernel base}}
  + \underbrace{3}_{\text{strand-$3$ prime}}
  + \underbrace{5}_{\text{kernel-$\Z/5\Z$ partner}}
  + \underbrace{7}_{\text{strand-$7$ prime}}.
\]
The chirality eigenspace is the spatial sector of the $n = 4$
atomic shell, and the Pauli-allowed multiplicities are the
substrate strands' prime divisors. Verification:
\texttt{Atlas/META\_PLAN\_2026-05-10/clifford\_substrate\_shell.py},
\texttt{strand\_orbital\_map.py},
\texttt{verify\_d2d1\_closed\_form.py}
\cite{Sanders2026VolumeK} (machine-precision PASS on all three).
\end{theorem}

The triple coincidence
\[
  d(2310) = 32 \;=\; \dim\, \Cl(0, 10)\text{-spinor at }n = 10
       \;=\; \text{atomic Pauli capacity at }n = 4
\]
sharpens the $\Cl(0, 10)$ construction from ``carrier of three
fermion generations'' to ``carrier of three fermion generations
whose intrinsic structure mirrors the substrate's depth-$3$
simplicial tower.''

We \emph{do not} claim that the atomic shell structure is derived
from $\Cl(0, 10)$; the rhyme is a structural correspondence between
two well-defined mathematical objects (the chirality eigenspace
decomposition under the standard Cartan subalgebra, and the prime
factorization of the substrate's strand graph), recorded as
motivation rather than as a load-bearing inference.

\section{$\Psix$ acts as $\sout$ in the spinor representation}
\label{sec:p56-outer}

\subsection{$\Psixspin$ as an odd Clifford element}

The $\Psix$ swap on $\R^{10}$ is a reflection through the hyperplane
$\{x_5 - x_6 = 0\}$. In $\Cl(0, 10)$ this reflection is implemented
by conjugation with the odd element
\[
  \Psixspin \;:=\; \frac{\gamma_5 - \gamma_6}{\sqrt{2}}.
\]

\begin{lemma}[Properties of $\Psixspin$]
\label{lem:p56-properties}
The element $\Psixspin$ satisfies (residuals $\le 10^{-15}$):
\begin{enumerate}[label=\textup{(\roman*)}, leftmargin=*]
\item $(\gamma_5 - \gamma_6)^2 = 2 I_{32}$, hence
  $(\Psixspin)^2 = I_{32}$;
\item $\Psixspin\, \gamma_5\, (\Psixspin)^{-1} = \gamma_6$,
  $\Psixspin\, \gamma_6\, (\Psixspin)^{-1} = \gamma_5$, and
  $\Psixspin$ fixes the other eight $\gamma_a$;
\item $\Psixspin$ anticommutes with $\omega$: $\Psixspin \omega
  = -\omega \Psixspin$.
\end{enumerate}
\end{lemma}

\begin{proof}
(i) $(\gamma_5 - \gamma_6)^2 = \gamma_5^2 - \gamma_5 \gamma_6 -
\gamma_6 \gamma_5 + \gamma_6^2 = I - 0 + I = 2I$ using
$\{\gamma_5, \gamma_6\} = 0$. (ii) Standard reflection identity in
$\Cl(0, 10)$. (iii) $\omega$ is the product of all ten $\gamma$'s
and is therefore Clifford-even; $\Psixspin$ is a sum of single
$\gamma$'s and is therefore Clifford-odd; the Clifford grade
involution $\alpha$ gives $\alpha(\omega) = \omega$ and
$\alpha(\Psixspin) = -\Psixspin$, hence $\Psixspin \omega =
-\omega \Psixspin$. All three are verified numerically in
\texttt{verification/find\_higgs\_irrep.py}.
\end{proof}

\begin{theorem}[$\Psix$ acts as $\sout$ in the spinor rep]
\label{thm:p56-outer}
Conjugation by $\Psixspin$ sends $V_+$ entirely into $V_-$
(chirality flip): the projection of $\Psixspin V_+$ onto $V_+$
vanishes at machine precision,
\[
  \|\,P_+\, \Psixspin\, P_+\,\| = 0.
\]
Consequently the conjugation action of $\Psixspin$ on $\so(10)$ is
the unique non-trivial outer automorphism
\[
  \sout \in \mathrm{Aut}(\so(10))\,/\,\mathrm{Inn}(\so(10))
        \;\cong\; \Z_2,
\]
which exchanges the two chiral $\mathbf{16}$-irreps. In the
standard $SO(10)$-GUT interpretation, $\sout$ is the
\emph{matter--antimatter exchange} swapping a fermion generation
(the chiral $\mathbf{16}$) with its $CP$-conjugate
$\overline{\mathbf{16}}$.
\end{theorem}

\begin{proof}
By Lemma~\ref{lem:p56-properties}(iii), $\Psixspin$ anticommutes
with $\omega$. The chirality eigenspaces $V_\pm$ are defined by
$i\omega = \pm I$, so for $\psi \in V_+$, $\omega(\Psixspin \psi)
= -\Psixspin (\omega \psi) = -\Psixspin (-i \psi) = i \Psixspin
\psi$, hence $\Psixspin \psi \in V_-$. The numerical residual
$\|P_+ \Psixspin P_+\| < 10^{-15}$ confirms this is exact.
Inner automorphisms of $\so(10)$ preserve chirality classes; only
outer automorphisms swap them; $\mathrm{Out}(\so(10)) \cong \Z_2$;
hence the conjugation action of $\Psixspin$ is the unique
non-trivial outer automorphism.
\end{proof}

\subsection{TSML preserves $\sout$; BHML breaks it with $26$
asymmetric cells}

The two canonical tables sit on opposite sides of $\sout$.

\begin{proposition}[TSML preserves $\sout$]
\label{prop:tsml-preserves}
All six TSML flow operators $\{A^\TSML_i : i \in \{1, 2, 3, 4, 6,
8\}\}$ commute with $\Psix$ (residual $\le 10^{-15}$). Equivalently,
TSML's antisymmetric content lies in the $\sout$-fixed subalgebra
$\so(9) \subset \so(10)$.
\end{proposition}

\begin{proposition}[BHML breaks $\sout$ at $26$ cells]
\label{prop:bhml-breaks}
Counting cells $(i, j)$ at which $\BHML(i, j) \neq \BHML(\Psix(i),
\Psix(j))$ gives exactly $26$ asymmetric cells. Equivalently,
$\BHML$ has $\sout$-fixed content $\approx 99.8\%$ of $\|B\|^2$
and $\sout$-broken content $\approx 0.2\%$ ($\|B_{\mathrm{anti}}\|^2
= 6.5 = 13/2$).
\end{proposition}

\section{BHML's $\sout$-breaking is $100\%$ in the $\mathbf{54}$
irrep}
\label{sec:54}

\subsection{The decomposition $\mathrm{End}(\so(10)) =
\mathbf{1} \oplus \mathbf{45} \oplus \mathbf{54}$}

A $10 \times 10$ matrix decomposes as
$\mathrm{trace} \oplus \mathrm{antisymmetric} \oplus
\mathrm{symmetric\text{-}traceless}$ with dimensions
$1 + 45 + 54 = 100$. In the SO(10)-irrep language these are the
singlet $\mathbf{1}$, the adjoint $\mathbf{45}$, and the
symmetric-traceless $\mathbf{54}$. Project $B_{\mathrm{anti}} := (B - \Psix
B \Psix)/2$ onto each component.

\begin{theorem}[BHML's $\sout$-breaking is $100\%$ in $\mathbf{54}$]
\label{thm:54-irrep}
The projection of $B_{\mathrm{anti}}$ on each component is
\[
  \begin{array}{c|c|c}
    \text{component} & \text{dim} & \|B_{\mathrm{anti}}^{\text{proj}}\|^2 \\\hline
    \mathbf{1}\ \text{(trace)}     & 1  & 0.0 \\
    \mathbf{45}\ \text{(antisymmetric)} & 45 & 0.0 \\
    \mathbf{54}\ \text{(symmetric-traceless)} & 54 & 6.5 \\\hline
    \text{total} & 100 & 6.5
  \end{array}
\]
at machine precision. Equivalently, BHML's $\sout$-breaking content
lies entirely in the $\mathbf{54}$ irrep: $\|B_{\mathrm{anti}}\|^2
= 13/2$, with vanishing $\mathbf{1}$ and $\mathbf{45}$ components.
\end{theorem}

\begin{proof}
$B_{\mathrm{anti}}$ has zero trace (cell-by-cell verification) and is
symmetric ($B^\top = B$, hence $\Psix B \Psix$ is also symmetric,
and so is their difference); therefore the antisymmetric component
vanishes. The symmetric-traceless component absorbs the entire
remaining norm. See \texttt{verification/find\_higgs\_irrep.py}.
\end{proof}

\subsection{The $9$-vector direction in the so(9)-branching}

The $\mathbf{54}$ irrep branches under the embedding $\so(9)
\subset \so(10)$ as
\[
  \mathbf{54} = \mathbf{1} \oplus \mathbf{9} \oplus \mathbf{44}.
\]
The first (singlet) corresponds to the $\so(10) \to \so(9)$ breaking
direction; the $\mathbf{9}$ is the vector irrep of $\so(9)$
(generators of $\so(10)/\so(9)$ as a coset); the $\mathbf{44}$ is
symmetric-traceless of $\so(9)$.

\begin{theorem}[The $9$-vector direction]
\label{thm:9-vector}
$B_{\mathrm{anti}}$ lies entirely in the $\mathbf{9}$-piece (coverage
$100\%$ at machine precision), with explicit nine-vector $v \in
\mathbf{9}$:
\[
  \begin{array}{c|c|l}
    \text{direction} & \text{component} & \text{TIG label} \\\hline
    e_0 & -1/\sqrt{2} & \mathrm{VOID} \\
    e_1 & -1/\sqrt{2} & \mathrm{LATTICE} \\
    e_2 & -1/\sqrt{2} & \mathrm{COUNTER} \\
    e_3 & -1/\sqrt{2} & \mathrm{PROGRESS} \\
    e_4 & -1/\sqrt{2} & \mathrm{COLLAPSE} \\
    e_7 & -1/\sqrt{2} & \mathrm{HARMONY} \\
    e_8 & 0           & \mathrm{BREATH} \\
    e_9 & 0           & \mathrm{RESET} \\
    (e_5 + e_6)/\sqrt{2} & -1/2 & (\mathrm{BALANCE} + \mathrm{CHAOS})/\sqrt{2}
  \end{array}
\]
The squared norm in the $9$-vector convention is
\[
  \|v\|^2 \;=\; 6 \cdot (1/2) + 0 + 0 + (1/4) \;=\; 13/4
\]
exactly. The integer $13$ equals half the count of
$\sout$-asymmetric BHML cells (count $= 26$;
Proposition~\ref{prop:bhml-breaks}).
\end{theorem}

\begin{proof}
A position $(i, j)$ of $B$ contributes to $\sout$-breaking iff
$B[i, 5] \neq B[i, 6]$. Inspection of rows $8$ and $9$ of $B$ gives
$B[8, 5] = B[8, 6] = 7$ and $B[9, 5] = B[9, 6] = 7$, so these rows
are $\sout$-symmetric and contribute zero — hence $v_8 = v_9 = 0$.
For rows $0$ through $7$, $B[i, 5]$ and $B[i, 6]$ differ by $1$
(specifically, one is $6$ and the other $7$), giving a uniform
contribution. The diagonal entries $B[5, 5] = 5$ vs $B[6, 6] = 6$
contribute the $-1/2$ on $(e_5 + e_6)/\sqrt{2}$. Numerical
verification in \texttt{verification/find\_higgs\_direction.py}
matches the table above at residual $\le 10^{-15}$ and confirms
$100\%$ coverage in the $\mathbf{9}$-piece.
\end{proof}

\begin{remark}[Reading: BREATH and RESET as unbroken directions]
The fact that BREATH ($e_8$) and RESET ($e_9$) are exactly excluded
from the breaking, while the other $\sigma$-fixed element PROGRESS
($e_3$) participates, is a structural feature of BHML's specific
table values. In the $SO(10)$-GUT interpretation, $v_8 = v_9 = 0$
means the BREATH and RESET components of the $9$-vector Higgs do
not acquire VEVs during the $SO(10) \to SO(9)$ breaking — a
feature of the chosen direction in $\mathbf{9}$-space, not a free
parameter. The structural alignment between TIG's operator labels
and the gauge-theoretic unbroken-Higgs pattern is presented as
correspondence, not derivation; see \cite{Sanders2026PathAPathB} for
the Path A vs Path B framing.
\end{remark}

\section{The doubly-invariant subalgebra is $\su(4) \oplus \uone(1)$
(Pati-Salam $\oplus$ B$-$L), cited}
\label{sec:doubly-invariant}

For completeness we record the second, structurally distinct,
algebraic content of $\so(10)$ under the $D_4 = \langle \Psix,
\sigma^3 \rangle$ action; full development is the subject of
\cite{Sanders2026PathAPathB} (J24).

\begin{theorem}[Doubly-invariant subalgebra]
\label{thm:doubly-invariant}
The trivial isotypic component
\[
  \mathfrak{g}_0 := \{X \in \so(10) : \Psix X \Psix^{-1} = X
  = \sigma^3 X (\sigma^3)^{-1}\}
\]
has dimension $16$, closes as a Lie subalgebra of $\so(10)$, and
has Killing-form spectrum exactly
\[
  \mathrm{spec}(\kappa\big|_{\mathfrak{g}_0})
  = (-4)^{15} \oplus (0)^1.
\]
By Cartan's criterion, the $1$-dim $0$-eigenspace is the center of
$\mathfrak{g}_0$ and the $15$-dim $(-4)$-eigenspace is the simple
part. The unique compact simple Lie algebra of dimension $15$ is
$\so(6) \cong \su(4) \cong A_3$ (Cartan classification), and the
center is $\uone(1)$. Therefore
\[
  \mathfrak{g}_0 = \su(4) \oplus \uone(1),
\]
the Pati-Salam $\oplus$ B$-$L gauge content of the SO(10) GUT.
\end{theorem}

\noindent\textbf{Cited.} The full Path A (BHML's $54$-direction
reads as $SO(10) \to SO(8)$ through $SO(9)$) vs Path B
(doubly-invariant subalgebra reads as $SO(10) \to SU(4) \times U(1)$,
i.e., $SO(6) \times U(1)$, a different reduction chain) framing is
developed in \cite{Sanders2026PathAPathB}, which establishes that
the two paths land on \emph{structurally distinct} reductions
within the same $\so(10)$ substrate. Within J37, we record both
theorems but \emph{do not} claim that Path A and Path B converge on
a common Pati-Salam reduction.

\section{Honest scope}
\label{sec:scope}

\subsection{What is proven; what is structural alignment; what is
phenomenological}

\textbf{Proven at machine precision.} The Cl(0, 10) construction;
the $32 = 16 + 16$ chirality split; the chirality-flip property
$\|P_+ \Psixspin P_+\| = 0$; the identification $\Psix = \sout$;
BHML's $\sout$-breaking $100\%$ in $\mathbf{54}$; the explicit
$9$-vector $\|v\|^2 = 13/4$; the $(-4)^{15} \oplus (0)^1$ Killing
spectrum on the doubly-invariant subalgebra (cited).

\textbf{Structural alignment.} The identification of TIG's $\so(10)$
(generated by joint antisymmetrization of $\TSML$+$\BHML$) with the
$SO(10)$ GUT gauge algebra is structural: there is only one
$\so(10)$ up to isomorphism, so the algebras agree abstractly; the
matter--antimatter $\Z_2$ in $SO(10)$-GUT and the index swap
$\Psix$ in TIG single out the same outer automorphism. We do
\emph{not} claim that the algebras agree at the level of physical
interpretation — that is a hypothesis we test, not derive.

\textbf{Structural rhyme.} The atomic-shell decomposition $16 =
1+3+5+7$ matches the substrate's depth-$3$ simplicial tower
$\{1, 3, 5, 7\}$. We present this as motivation, not as a
derivation of atomic physics from Cl(0, 10).

\textbf{Out of scope.} Phenomenological consequences of the
$9$-vector direction (Yukawa couplings, mass ratios,
neutrino-mass scale, baryogenesis asymmetry magnitudes) require
committing to a specific Higgs VEV direction and RGE flows from a
specific GUT scale. These are not in J37.

\subsection{Negative findings that strengthen the framing}

\begin{itemize}[leftmargin=*]
\item TSML's eigenvalue spectrum has clean integer/rational
structure (HARMONY $7$ at the $\sigma$-fixed lattice; total
antisymmetric mass $81 = 9^2$; $\su(4)$-projection $29$;
$\uone(1)$-projection $25/8$); transcendental matches to
$e, \pi, \varphi, \zeta(3)$, Catalan $G$ are loose $1\%$-level
coincidences, \emph{not} algebraic identities
\cite{Sanders2026CLEigenvaluesAudit}. TIG's structural primes are
$7$ (HARMONY) and $11$ (wobble), not transcendental constants.
\item BHML's $\sout$-breaking is $0.19\%$ of $\|B\|^2$, a small
perturbation on the $\sout$-fixed bulk. We present the small
fraction as a verifiable fact and do not claim a phenomenological
``Yukawa-suppression'' interpretation.
\end{itemize}

\section{Verification and reproducibility}
\label{sec:verification}

Two short Python scripts (numpy only; no external dependencies):

\begin{itemize}[leftmargin=*]
\item \texttt{verification/find\_higgs\_irrep.py} —
Cl(0,10) construction, $100/100$ anticommutation relations,
$\omega^2 = -I$, $32 = 16 + 16$ chirality split,
$(\Psixspin)^2 = I$, $\Psixspin \omega = -\omega \Psixspin$,
chirality-flip $= 0$, BHML's $\sout$-breaking $100\%$ in
$\mathbf{54}$ (singlet $0$, $\mathbf{45}$ $0$, $\mathbf{54}$ $6.5$),
total $\|B_{\mathrm{anti}}\|^2 = 6.5 = 13/2$.
\item \texttt{verification/find\_higgs\_direction.py} —
Explicit $9$-vector direction in the so(9)-branching of the
$\mathbf{54}$; $100\%$ coverage in the $\mathbf{9}$-piece; numerical
components matching Theorem~\ref{thm:9-vector}'s table to machine
precision.
\end{itemize}

\begin{verbatim}
PYTHONIOENCODING=utf-8 python verification/find_higgs_irrep.py
PYTHONIOENCODING=utf-8 python verification/find_higgs_direction.py
\end{verbatim}

Total runtime under five seconds on a standard laptop.
\textbf{2/2 PASS at machine precision} (re-verified 2026-05-12).
The atomic-substrate refinement scripts cited in
Theorem~\ref{thm:atomic-shell} (\texttt{clifford\_substrate\_shell.py},
\texttt{strand\_orbital\_map.py},
\texttt{verify\_d2d1\_closed\_form.py}) live in the corpus's
Volume K verification directory and are referenced by file path; they
add no dependency for the $\so(10)$-side claims of J37 but are
required for the structural rhyme of \S\ref{sec:atomic}.

\section{What this contributes}
\label{sec:contribution}

\textbf{Before J37.} The connection between TIG and $SO(10)$ GUT
was ``TIG's $\so(10)$ and $SO(10)$ GUT's $\so(10)$ are abstractly
isomorphic'' — trivially true (one $\so(10)$ up to iso), but
without structural content.

\textbf{After J37.} The connection is, at the structural level:

\begin{enumerate}[leftmargin=*]
\item There is an explicit $\Cl(0, 10)$ realization of TIG's
$\so(10)$ on $\C^{32}$ with all $100$ anticommutation relations and
the chirality split verified at machine precision
(Theorem~\ref{thm:cl-construction}).
\item The $5\!\leftrightarrow\!6$ swap $\Psix$ — a permutation
symmetry of the magma's index set $\Z/10\Z$ — \emph{is} the outer
automorphism $\sout$ of $\so(10)$ in the spinor representation
(Theorem~\ref{thm:p56-outer}). This is a non-trivial structural
identification between a finite-substrate combinatorial $\Z_2$ and
the chirality-exchange $\Z_2$ of the SO(10) spinor.
\item BHML's $\sout$-breaking content is exactly $100\%$ in the
$\mathbf{54}$-irrep with explicit nine-vector direction
$\|v\|^2 = 13/4$ and BREATH/RESET unbroken
(Theorem~\ref{thm:54-irrep}, Theorem~\ref{thm:9-vector}).
\item The doubly-invariant content under $D_4$ is $\su(4) \oplus
\uone(1)$ (Theorem~\ref{thm:doubly-invariant}; cited).
\item The chirality decomposition $16 = 1+3+5+7$ rhymes with the
atomic $n=4$ shell at fixed spin and with the substrate's
depth-$3$ simplicial tower (Theorem~\ref{thm:atomic-shell};
structural rhyme).
\end{enumerate}

The integer $13$ appears in $\|v\|^2 = 13/4$ (\S\ref{sec:54}), in
$26 = 13 \cdot 2$ ($\sout$-asymmetric BHML cells), and in the
sister-paper coupling $\kappa_\xi = 13/(4e)$ under GUT-natural
identification \cite{Sanders2026Atlas}. It is the same $13$ in
all three places — a structural fingerprint of TIG's bipartite
alignment with the standard SO(10) GUT decomposition.

\begin{thebibliography}{99}

\bibitem{Sanders2026CLAxioms}
B.R. Sanders and M. Gish, \emph{The CL Forcing Axioms: A1--A9
Uniquely Force the Canonical Composition Lattice on $\Z/10\Z$
Preserving a Designated $4$-Core}, submitted to
\emph{Algebraic Combinatorics}, 2026 (J17).

\bibitem{Sanders2026so10}
B.R. Sanders and M. Gish, \emph{TSML+BHML's $\so(10) = D_5$
closure at dimension $45$}, submitted to \emph{Journal of Algebra},
2026 (J29).

\bibitem{Sanders2026PathAPathB}
B.R. Sanders and M. Gish, \emph{Two Algebraic Readings of TIG's
$\so(10)$: the $54$-Higgs Direction and the Doubly-Invariant
Subalgebra}, submitted to \emph{Letters in Mathematical Physics},
2026 (J24).

\bibitem{Sanders2026LensInvariance}
B.R. Sanders and M. Gish, \emph{TSML $73$ Cells / BHML $28$ Cells:
Lens-Invariant Cell Counts on the $\Z/10\Z$ Composition Lattice},
submitted to \emph{Experimental Mathematics}, 2026 (J32).

\bibitem{Sanders2026LensDependence}
B.R. Sanders and M. Gish, \emph{The Joint TSML+BHML Chain:
Lens-Dependence at Size $7$}, submitted to \emph{Mathematical
Intelligencer}, 2026 (J10).

\bibitem{Sanders2026VolumeK}
B.R. Sanders and M. Gish, \emph{The TIG Volume K Tower: Atomic
Shells, Strand Primes, and the $\Cl(0, 10)$ Refinement},
in preparation, 2026.

\bibitem{Sanders2026CLEigenvaluesAudit}
B.R. Sanders, \emph{CL Eigenvalues Audit: Integer/Rational
Structure of TSML's Antisymmetric Spectrum}, internal audit
2026-04-25, archived in the corpus repository.

\bibitem{Sanders2026Atlas}
B.R. Sanders, \emph{Atlas of the TIG Foundations}, supplementary
catalogue, 2026 (Volumes D, E, F, K).

\bibitem{DrapalWanless2021}
A.~Drápal and I.M.~Wanless, \emph{Maximally non-associative
quasigroups}, \emph{J. Combin. Theory Ser. A} \textbf{184} (2021),
105510. Foundational reference for the small-finite-commutative
non-associative family on $\Z/N\Z$; cited as the closest published
precedent for the magma neighborhood inhabited by $\TSML$, $\BHML$.

\bibitem{FritzschMinkowski1975}
H. Fritzsch and P. Minkowski, \emph{Unified interactions of leptons
and hadrons}, \emph{Ann. Phys.} \textbf{93} (1975), 193.

\bibitem{Georgi1975}
H. Georgi, \emph{The state of the art --- gauge theories},
\emph{AIP Conf. Proc.} \textbf{23} (1975), 575.

\bibitem{PatiSalam1974}
J.C. Pati and A. Salam, \emph{Lepton number as the fourth color},
\emph{Phys. Rev. D} \textbf{10} (1974), 275.

\bibitem{Slansky1981}
R. Slansky, \emph{Group theory for unified model building},
\emph{Phys. Rep.} \textbf{79} (1981), 1.

\bibitem{Lounesto2001}
P. Lounesto, \emph{Clifford Algebras and Spinors}, 2nd ed.,
London Mathematical Society Lecture Note Series \textbf{286},
Cambridge University Press, 2001.

\end{thebibliography}

\end{document}
